\documentclass[../main.tex]{subfiles}

\begin{document}

\chapter{Aturan Inferensi dan Penarikan Kesimpulan}

\begin{subcpmk}
  \item Mahasiswa mampu menganalisis validitas argumen menggunakan aturan inferensi deduktif matematis murni (logika silogistik).
  \item Mahasiswa mengenali pola Modus Ponens, Modus Tollens, dan skema Silogisme Hipotetis dalam memecahkan persoalan rasional.
\end{subcpmk}

\section{Pendahuluan Penarikan Simpulan Komputatif}

Sistem pengambil keputusan \textit{(Decision Support Systems)} dalam perangkat lunak modern digerakkan oleh *Inference Engines*. Di belakang *Inference Engines* (Mesin Inferensi komputasi *Artificial Intelligence* level awalan), bersemayam kemampuan dan kerangka untuk mengidentifikasi sekumpulan *Premis* aksiomatik maupun deklarasi faktual untuk, \textit{melalui penalaran silogistik deduktif presisi}, menarik dan memutuskan turunan berupa satu *Konklusi (Kesimpulan)* berdaya valid dan \textit{reliable} kuat.

Aturan Inferensi (\textit{Rules of Inference}) memandu bagaimana dari kondisi-kondisi (premis-premis matematis) $P_1, P_2, P_3 \dots P_k$ berimbas ke Konklusi logis berdinamika konsisten ($Q$).

Suatu inferensi (argumen) dijuluki \textbf{VALID} takkala, jika \textit{seluruh susunan dari premis-premis penggeraknya bernilai aktual Benar (T)}, maka keniscayaan akan memastikan Konklusi-nya wajib Bernilai \textbf{Benar (T)}. Pembuktian teknis ini direduksi setara ke ekivalensi dari Tautologi majemuk $(P_1 \land P_2 \land \dots \land P_k) \rightarrow Q$.

\section{Aturan Pokok Inferensi Proposisional Dasar}

\subsection{Modus Ponens (Hukum Pembenaran Khusus)}
Juga termasyhur selaku \textit{The Rule of Detachment}. Esensinya menyatakan bilamana sebuah prasyarat kebenaran komputasional kausal dipenuhi premis bersyarat, dan pre-kondisi itu pun benar-benar konkret/terjadi mutlak, maka kesimpulan dan \textit{trigger} akibat bersyaratnya logis secara otomatis dieksekusi benar. Ekstraknya:\\
Premis Mayor: $p \rightarrow q$ (Jika hujan interupsi, maka \textit{trigger sensor air} aktif) \\
Premis Minor: $p$ (Memang turunlah curah hujan iterupsi) \\
Kesimpulan: $\therefore q$ (Sirkuit menyatakan \textit{sensor air divalidasi aktif!})

\subsection{Modus Tollens (Hukum Penolakan)}
Menyitir dari Modus Ponens, namun diinisiasikan di sisi pengingkaran akibat semantik (\textit{The Law of Contrapositive Denial}). \\
Premis Mayor: $p \rightarrow q$ \\
Premis Minor: $\sim q$ (Diinspeksi ternyata: \textit{sensor air bernilai matifalse / False!}) \\
Kesimpulan: $\therefore \sim p$ (Otomatis ditarik validitas inferensi komputasi Logis konklusinya: Saat ini Hujan Interupsi jelas tak berlaku = False / Negasi $p$).

\subsection{Silogisme Hipotetis}
Bila mesin algoritma AI disuplai rangkaian berantai kognisi berdimensi kontinuitas dari $p$ meluncur ke $q$, dibarengi $q$ bermuara meluncur pada $r$, logis mereduksi kesimpulan transisional silogisme (hukum transisi).\\
Premis 1: $p \rightarrow q$ \\
Premis 2: $q \rightarrow r$ \\
Kesimpulan: $\therefore p \rightarrow r$

\subsection{Silogisme Disjungtif}
Menentukan keputusan dari gerbang komputasi persilangan OR (\textit{The Law of Disjunctive Inference}).\\
Premis 1: $p \lor q$ (Sistem \textit{booting} memakai Partisi X Atau Sistem Partisi Induk Utama/Y)\\
Premis 2: $\sim p$ (Log \textit{compiler} memindai tak ada \textit{booting} lewat partisi dasar X)\\
Kesimpulan: $\therefore q$ (Secara deduksi logis, PC ini beroperasi \textit{booting} spesifik via Partisi Induk Y!)

\begin{aktivitas}
  \item \textbf{Diskusi Rasional AI Base}: Sistem navigasi taktis menyatakan \textit{(A1)} Jika bensin habis maka fungsi motor utama putus. \textit{(A2)} Ternyata saat ini fungsi motor utama tidak terputus beroperasi. Apa inferensi mutlak status bahan bensin Anda dan memakai logika aturan Aturan Pokok Inferensi manakah itu?
\end{aktivitas}

\begin{latihan}
  \item Terjemahkan pola inferensial ini, tentukan ia tergolong Valid argumen dari aturan spesifik yang mana:
  "Jika hari Senin, \textit{database cronjob backup} dieksekusi secara otomatis." (Premis Mayor).
  "Dan juga \textit{Backup database cronjob} telah terevaluasi dan batal/tak berjalan di server (False)." (Premis Minor).
  Kesimpulan : "Bisa disimpulkan dengan yakin formal algoritmanya hari jadwal saat ini *Bukan Senin*"
\end{latihan}

\begin{checklist}
  \item Mampu menspesifikan parameter deduktif pada pengkondisian skrip fungsional yang valid.
  \item Hafal 4 Aturan Pokok: Ponens, Tollens, Transisi Silogisme Hipotetis dan Disjungtif Logis.
\end{checklist}

\end{document}
